\section{Knot Invariants}

\subsection{Alexander Polynomial}
A knot invariant assigns the same value to every diagram
of the same knot. One important example is the Alexander polynomial $\Delta_K(t)$, defined up to multiplication by $\pm t^m$
m. A common normalization satisfies

\begin{equation}
    \Delta_k(t) = \Delta_K(t^{-1}, \quad \Delta_k(1)=1
\end{equation}

For the trefoil knot,
\begin{equation}
    \Delta_{3_1} (t) = t^{-1} -1 + t
\end{equation}
The determinant of a knot is obtained from the polynomial by
\begin{equation}
    \det (K) = | \Delta_K(-1)|.
\end{equation}
Hence,
\begin{equation*}
    \det(3_1) = | -1-1-1| = 3
\end{equation*}
Table~\ref{tab:inv} compares these quantities for several small
knots.


\subsection{A Skein Relation}
Many polynomial invariants are computed by comparing
three diagrams that differ only inside a small region. For
the Alexander polynomial, one form of the skein relation
is
\begin{equation}
    \Delta_{L_+} (t) - \Delta_{L_-} (t) = \left( t^{1/2} - t^{-1/2}\right) \Delta_{L_0}(t)
\end{equation}

where $ L_+, L_-$, and $ L_0$ denote a positive crossing, a negative crossing, and a smoothing, respectively.




\subsection{The Signature of a Knot}
Another useful knot invariant is the signature $\sigma(K)$, obtained from a Seifert matrix $V$ . It is defined as the signature of the symmetric matrix
\[
V+V^T,
\]

that is, the number of positive eigenvalues minus the
number of negative eigenvalues.

For a knot and its mirror image, the signature satisfies
\begin{equation}
    \sigma(K) = \begin{cases}
        0, & \text{if $K$ is the unknot,} \\
        -2, & \text{is K is the right-handed trefoil}\\
        2, & \text{if K is the left-handed trefoil}
    \end{cases}
\end{equation}